% Modul Belajar 6, dihasilkan oleh tools/build.py dari tools/konten/p06.py
\documentclass[a4paper]{article}
\usepackage{modulITERA}
\begin{document}
\Sampul{6}{Perulangan Bagian 2 dan Pola}{Perulangan bersarang, break dan continue, serta menurunkan rumus pola}{CPMK0607 (menerapkan) dan CPMK0608 (menjelaskan) pada konsep dasar algoritma dan sistem komputer}{sekitar 3 jam belajar mandiri, ditambah 150 menit di kelas}{\texttt{02 Hands-on/P6 - Perulangan 2 dan Pola - Hands-on/}}
\Bagian{Modul 6}{Peta belajar}
Ikuti urutan ini dari atas ke bawah. Setiap bagian memakai apa yang sudah dibahas di bagian sebelumnya.
\par\vspace{4pt}\noindent{\fontsize{9.5pt}{13pt}\selectfont
\begin{tabular}{@{}>{\raggedright\arraybackslash}p{0.140\textwidth}>{\raggedright\arraybackslash}p{0.840\textwidth}@{}}
\kepala{Urutan} & \kepala{Bagian}\\
1 & Tentang modul ini\\\garis
2 & Masalah pembuka: Denah kursi bioskop\\\garis
3 & Langkah 1: Perulangan bersarang\\\garis
4 & Langkah 2: Batas dalam bergantung pada baris\\\garis
5 & Langkah 3: Menurunkan rumus pola\\\garis
6 & Langkah 4: Piramida dari rumus\\\garis
7 & Langkah 5: break dan continue\\\garis
8 & Langkah 6: Mencari dengan break\\\garis
9 & Pesan error yang sering muncul\\\garis
10 & Latihan mandiri\\\garis
11 & Rangkuman\\\garis
12 & Cek pemahaman\\\garis
13 & Exit Ticket dan tugas\\\garis
14 & Bacaan lanjut dan persiapan minggu depan\\
\garisdasar
\end{tabular}}\par\vspace{6pt}
\Bagian{Pembuka}{Tentang modul ini}
Modul ini mendampingi Pertemuan 6, pertemuan terakhir sebelum review UTS. Kalian akan menaruh perulangan di dalam perulangan, yang membuka jalan untuk pola, tabel, dan nanti array dua dimensi. Kunci minggu ini bukan menghafal kode pola, tetapi menurunkan rumusnya.

Isinya sama dengan slide di kelas, tetapi ditulis lebih pelan dan lebih lengkap, supaya kalian bisa mengulang sendiri di rumah tanpa harus menebak maksud slide.

\Sub{Setelah menyelesaikan modul ini, kalian bisa}
\par\vspace{4pt}\noindent{\fontsize{9.5pt}{13pt}\selectfont
\begin{tabular}{@{}>{\raggedright\arraybackslash}p{0.070\textwidth}>{\raggedright\arraybackslash}p{0.680\textwidth}>{\raggedright\arraybackslash}p{0.230\textwidth}@{}}
\kepala{No} & \kepala{Kemampuan} & \kepala{Dibahas di}\\
1 & Menulis program dengan perulangan bersarang, \texttt{break}, dan \texttt{continue} untuk menghasilkan pola dan memproses data dua dimensi (Sub-CPMK 6.1, CPMK0607) & Langkah 1 sampai 6\\\garis
2 & Menurunkan hubungan antara nomor baris dan banyak isi pada sebuah pola, serta menelusuri perulangan bersarang untuk menentukan keluarannya (Sub-CPMK 6.2, CPMK0608) & kotak Tebak dulu, Latihan 3\\
\garisdasar
\end{tabular}}\par\vspace{6pt}
Karena setiap instrumen penilaian dibagi rata antara Bagian A (menulis kode) dan Bagian B (menelusuri dan menjelaskan), yang dinilai bukan hanya program kalian jalan, tetapi juga kalian bisa menjelaskan mengapa program itu jalan.

\Sub{Cara memakai modul ini}
Setiap langkah punya pola yang sama: penjelasan konsep, kadang contoh dari kehidupan sehari-hari, lalu kode yang bisa langsung kalian ketik. Sesudah kode ada kotak \textbf{Tebak dulu}. Tulis tebakan kalian di kertas sebelum membaca \textbf{Jawaban} di bawahnya. Kebiasaan menebak lalu membuktikan inilah yang membuat konsep menempel, bukan sekadar membaca.

\begin{Penting}[Ketik, jangan salin]
Ketik ulang setiap contoh di GDB Online atau editor kalian, jangan disalin-tempel. Saat mengetik, kalian akan membuat salah ketik, membaca pesan error, dan memperbaikinya. Itu bagian dari belajar.
\end{Penting}
\Sub{Yang perlu disiapkan}
Peramban dengan akses ke onlinegdb.com, atau g++ di komputer kalian sendiri. Kertas dan alat tulis untuk menebak keluaran dan membuat trace table. Folder hands-on pertemuan ini dari repo mata kuliah.

\Bagian{Pembuka}{Masalah pembuka: Denah kursi bioskop}
Bioskop ingin menampilkan denah kursi: 5 baris, tiap baris 8 kursi, dengan nomor seperti A1 sampai E8.

Satu perulangan cukup untuk satu baris kursi. Untuk semua baris, perulangan itu harus diulang lagi.

\begin{MKeluaran}[title={Tampilan yang diinginkan}]
A1 A2 A3 A4 A5 A6 A7 A8
B1 B2 B3 B4 B5 B6 B7 B8
C1 C2 C3 C4 C5 C6 C7 C8
D1 D2 D3 D4 D5 D6 D7 D8
E1 E2 E3 E4 E5 E6 E7 E8
\end{MKeluaran}
\begin{Tebak}[Coba pikirkan]
Mana yang berubah lebih cepat, huruf baris atau nomor kursi? Apa artinya bagi susunan perulangannya?
\end{Tebak}
\begin{Jawaban}[Cara memecahnya]
Pisahkan dimensi: Baris berganti lebih lambat, kolom berganti lebih cepat. Rumuskan: Untuk setiap baris, ulangi semua kolom; lalu pindah baris. Terjemahkan: \texttt{for} luar untuk baris, \texttt{for} dalam untuk kolom.
\end{Jawaban}
\Bagian{Langkah 1}{Perulangan bersarang}
Perulangan bersarang adalah perulangan di dalam badan perulangan lain. Untuk setiap satu putaran perulangan luar, perulangan dalam berjalan dari awal sampai selesai. Polanya mirip jam: jarum menit berputar penuh untuk setiap satu langkah jarum jam.

\texttt{char} juga bisa dipakai sebagai pencacah, karena setiap karakter punya kode angka yang berurutan.

\begin{MKode}{denah.cpp}
for (char b = 'A'; b <= 'C'; b++) {
    for (int k = 1; k <= 4; k++) {
        cout << b << k << " ";
    }
    cout << endl;
}
\end{MKode}
\begin{Tebak}[]
Sebelum menjalankan program, tulis keluarannya di kertas, karakter demi karakter.
\end{Tebak}
\begin{Jawaban}[]
\begin{MKeluaran}[]
A1 A2 A3 A4 
B1 B2 B3 B4 
C1 C2 C3 C4 
\end{MKeluaran}
Untuk setiap satu putaran \texttt{for} luar, \texttt{for} dalam berjalan penuh. Total putaran dalam: 3 × 4 = 12.
\end{Jawaban}
\Bagian{Langkah 2}{Batas dalam bergantung pada baris}
Pola terbentuk ketika batas perulangan dalam bergantung pada variabel perulangan luar. Pada segitiga, baris ke-b berisi tepat b bintang, sehingga batas dalamnya \texttt{k <= b}.

\begin{MKode}{segitiga.cpp}
int n = 4;
for (int b = 1; b <= n; b++) {
    for (int k = 1; k <= b; k++) {
        cout << "*";
    }
    cout << endl;
}
\end{MKode}
\begin{Tebak}[]
Sebelum menjalankan program, tulis keluarannya di kertas, karakter demi karakter.
\end{Tebak}
\begin{Jawaban}[]
\begin{MKeluaran}[]
*
**
***
****
\end{MKeluaran}
Batas \texttt{for} dalam memakai \texttt{b}. Itulah yang membuat setiap baris berbeda panjang.
\end{Jawaban}
\Bagian{Langkah 3}{Menurunkan rumus pola}
Cara paling andal membuat pola: tulis tabel nomor baris, banyak spasi, dan banyak isi untuk beberapa baris pertama. Dari tabel itu, rumus biasanya terlihat. Spasi berkurang satu per baris, jadi n − b. Bintang bertambah dua per baris mulai dari satu, jadi 2b − 1.

Setelah rumus ketemu, kodenya tinggal dua perulangan dalam yang berurutan: satu untuk spasi, satu untuk bintang.

\par\vspace{4pt}\noindent{\fontsize{9.5pt}{13pt}\selectfont
\begin{tabular}{@{}>{\raggedright\arraybackslash}p{0.152\textwidth}>{\raggedright\arraybackslash}p{0.174\textwidth}>{\raggedright\arraybackslash}p{0.174\textwidth}>{\raggedright\arraybackslash}p{0.479\textwidth}@{}}
\kepala{Baris b} & \kepala{Spasi} & \kepala{Bintang} & \kepala{Tampilan}\\
1 & 3 & 1 & \texttt{   *}\\\garis
2 & 2 & 3 & \texttt{  ***}\\\garis
3 & 1 & 5 & \texttt{ *****}\\\garis
4 & 0 & 7 & \texttt{*******}\\\garis
Rumus & n − b & 2b − 1 & dua \texttt{for} dalam berurutan\\
\garisdasar
\end{tabular}}\par\vspace{6pt}
\Bagian{Langkah 4}{Piramida dari rumus}
Kode berikut adalah terjemahan langsung dari tabel. Perhatikan bahwa batas setiap perulangan dalam persis sama dengan rumus di kolom tabel.

\begin{MKode}{piramida.cpp}
int n = 4;
for (int b = 1; b <= n; b++) {
    for (int s = 1; s <= n - b; s++) {
        cout << " ";
    }
    for (int k = 1; k <= 2 * b - 1; k++) {
        cout << "*";
    }
    cout << endl;
}
\end{MKode}
\begin{Tebak}[]
Sebelum menjalankan program, tulis keluarannya di kertas, karakter demi karakter.
\end{Tebak}
\begin{Jawaban}[]
\begin{MKeluaran}[]
   *
  ***
 *****
*******
\end{MKeluaran}

\end{Jawaban}
\Bagian{Langkah 5}{break dan continue}
\texttt{continue} melewati sisa badan dan langsung lanjut ke putaran berikutnya. \texttt{break} menghentikan perulangan sepenuhnya. Pada perulangan bersarang, keduanya hanya berlaku untuk perulangan terdalam yang memuatnya.

\begin{MKode}{breakcontinue.cpp}
for (int i = 1; i <= 10; i++) {
    if (i % 3 == 0) {
        continue;
    }
    if (i > 7) {
        break;
    }
    cout << i << " ";
}
cout << endl;
\end{MKode}
\begin{Tebak}[]
Sebelum menjalankan program, tulis keluarannya di kertas, karakter demi karakter.
\end{Tebak}
\begin{Jawaban}[]
\begin{MKeluaran}[]
1 2 4 5 7 
\end{MKeluaran}
\texttt{continue} melompat ke putaran berikutnya. \texttt{break} keluar dari perulangan terdekat saja, bukan semua perulangan.
\end{Jawaban}
\Bagian{Langkah 6}{Mencari dengan break}
\texttt{break} cocok untuk pencarian: begitu jawaban ditemukan, tidak ada gunanya terus memeriksa. Pada contoh ini, pemeriksaan pembagi berhenti di 7 karena 91 = 7 × 13. Variabel \texttt{bool} mencatat apakah pencarian berhasil.

\begin{MKode}{prima.cpp}
int x = 91;
bool prima = true;
for (int d = 2; d * d <= x; d++) {
    if (x % d == 0) {
        cout << "Habis dibagi " << d << endl;
        prima = false;
        break;
    }
}
cout << (prima ? "prima" : "bukan prima") << endl;
\end{MKode}
\begin{Tebak}[]
Sebelum menjalankan program, tulis keluarannya di kertas, karakter demi karakter.
\end{Tebak}
\begin{Jawaban}[]
\begin{MKeluaran}[]
Habis dibagi 7
bukan prima
\end{MKeluaran}

\end{Jawaban}
\Bagian{Waspada}{Pesan error yang sering muncul}
Semua pesan di tabel ini dihasilkan g++ dengan opsi \texttt{-Wall}, sama seperti yang dipakai \texttt{cek.sh}.
\par\vspace{4pt}\noindent{\fontsize{9.5pt}{13pt}\selectfont
\begin{tabular}{@{}>{\raggedright\arraybackslash}p{0.240\textwidth}>{\raggedright\arraybackslash}p{0.270\textwidth}>{\raggedright\arraybackslash}p{0.470\textwidth}@{}}
\kepala{Kesalahan} & \kepala{Contoh} & \kepala{Pesan pertama dari g++}\\
Variabel luar dideklarasi ulang & \texttt{for (int i...) for (int i...)} & \texttt{error: redeclaration of 'int i'}\\\garis
break di luar perulangan & \texttt{if (x) break;} & \texttt{error: break statement not within loop or switch}\\\garis
continue di luar perulangan & \texttt{continue; di main} & \texttt{error: continue statement not within a loop}\\\garis
Kurung kurawal kurang & \texttt{for \{ for \{ \}} & \texttt{error: expected '\}' at end of input}\\\garis
Pencacah dalam tak terpakai & \texttt{int baris; tak dipakai} & \texttt{warning: unused variable 'baris'}\\
\garisdasar
\end{tabular}}\par\vspace{6pt}
Kesalahan pola paling sering bukan error kompilasi, melainkan batas perulangan yang meleset satu. Tabel baris dan isi adalah alat utama untuk menemukannya.

\Bagian{Latihan}{Latihan mandiri}
Latihan ini sama dengan hands-on di kelas. Semua file ada di \texttt{02 Hands-on/P6 - Perulangan 2 dan Pola - Hands-on/latihan/}. Kerjakan berurutan, lalu cocokkan dengan \texttt{expected/} atau jalankan \texttt{bash cek.sh}.
\par\vspace{4pt}\noindent{\fontsize{9.5pt}{13pt}\selectfont
\begin{tabular}{@{}>{\raggedright\arraybackslash}p{0.070\textwidth}>{\raggedright\arraybackslash}p{0.360\textwidth}>{\raggedright\arraybackslash}p{0.150\textwidth}>{\raggedright\arraybackslash}p{0.400\textwidth}@{}}
\kepala{No} & \kepala{File} & \kepala{Tingkat} & \kepala{Yang dilatih}\\
1 & \texttt{handson1-persegi.cpp} & Dasar & \texttt{for} bersarang dengan batas tetap\\\garis
2 & \texttt{handson2-segitiga.cpp} & Menengah & batas dalam bergantung pada baris\\\garis
3 & \texttt{handson3-piramida.cpp} & Tantangan & tabel baris-isi, perbaiki tiga batas\\\garis
+ & \texttt{bonus-prima.cpp} & Pengayaan & bersarang dengan \texttt{break}\\
\garisdasar
\end{tabular}}\par\vspace{6pt}
\Sub{Latihan 1: handson1-persegi.cpp}
Persegi panjang bintang berukuran t × l. Masukan uji: \texttt{3 5}.

\Sub{Latihan 2: handson2-segitiga.cpp}
Segitiga angka dan segitiga bintang terbalik untuk n = 5.

\Sub{Latihan 3: handson3-piramida.cpp}
Piramida bintang yang batas perulangannya salah. Buat tabel baris, spasi, dan bintang; perbaiki; jelaskan.

\Sub{Bonus: bonus-prima.cpp}
Daftar bilangan prima sampai n dan banyaknya.

\Sub{Latihan menelusuri}
\begin{MKode}{tebak.cpp}
int hitung = 0;
for (int i = 1; i <= 3; i++) {
    for (int j = i; j <= 3; j++) {
        cout << i << j << " ";
        hitung++;
    }
}
cout << endl << hitung << endl;
\end{MKode}
\begin{Tebak}[]
Tulis keluarannya di kertas tanpa menjalankan program. Buat trace table bila perlu.
\end{Tebak}
\begin{Jawaban}[]
\begin{MKeluaran}[]
11 12 13 22 23 33 
6
\end{MKeluaran}
\texttt{j} mulai dari \texttt{i}, jadi baris luar ke-1 menghasilkan 3 pasangan, ke-2 dua, ke-3 satu. Total 6.
\end{Jawaban}
\Bagian{Penutup}{Rangkuman}
\par\vspace{4pt}\noindent{\fontsize{9.5pt}{13pt}\selectfont
\begin{tabular}{@{}>{\raggedright\arraybackslash}p{0.320\textwidth}>{\raggedright\arraybackslash}p{0.660\textwidth}@{}}
\kepala{Konsep} & \kepala{Yang perlu diingat}\\
Perulangan bersarang & Untuk setiap satu putaran \texttt{for} luar, \texttt{for} dalam berjalan penuh.\\\garis
Batas dalam bergantung pada baris & Batas \texttt{for} dalam memakai \texttt{b}.\\\garis
Menurunkan rumus pola & Cara paling andal membuat pola: tulis tabel nomor baris, banyak spasi, dan banyak isi untuk beberapa baris pertama.\\\garis
Piramida dari rumus & Kode berikut adalah terjemahan langsung dari tabel.\\\garis
break dan continue & \texttt{continue} melompat ke putaran berikutnya.\\\garis
Mencari dengan break & \texttt{break} cocok untuk pencarian: begitu jawaban ditemukan, tidak ada gunanya terus memeriksa.\\
\garisdasar
\end{tabular}}\par\vspace{6pt}
\Bagian{Penutup}{Cek pemahaman}
Jawab tanpa menjalankan program, lalu periksa dengan menjalankannya.
\begin{enumerate}
\item Berapa total \texttt{cout} yang dijalankan oleh dua \texttt{for} bersarang 4 × 6?
\item Pada segitiga terbalik tinggi 5, berapa bintang di baris ke-2?
\item Apa beda \texttt{break} dan \texttt{continue}?
\item Pada perulangan bersarang, \texttt{break} di perulangan dalam menghentikan apa?
\item Tuliskan rumus banyak spasi baris b pada piramida tinggi n.
\end{enumerate}
\begin{Jawaban}[]
\begin{enumerate}[leftmargin=16pt]
\item 24.
\item 4 (rumus n − b + 1).
\item \texttt{break} keluar dari perulangan; \texttt{continue} lompat ke putaran berikutnya.
\item Hanya perulangan dalam; perulangan luar tetap lanjut.
\item n − b.
\end{enumerate}
\end{Jawaban}
\Bagian{Penilaian}{Exit Ticket 6}
Dikerjakan di 25 menit terakhir kelas, sendiri, tanpa AI, internet, atau diskusi. Exit Ticket 6 adalah satu dari 14 Exit Ticket; rata-ratanya menjadi komponen berbobot 25\%.
\par\vspace{4pt}\noindent{\fontsize{9.5pt}{13pt}\selectfont
\begin{tabular}{@{}>{\raggedright\arraybackslash}p{0.080\textwidth}>{\raggedright\arraybackslash}p{0.600\textwidth}>{\raggedright\arraybackslash}p{0.100\textwidth}>{\raggedright\arraybackslash}p{0.200\textwidth}@{}}
\kepala{Soal} & \kepala{Isi} & \kepala{Poin} & \kepala{CPMK}\\
A1 & Tulis program yang menampilkan pola yang diberikan untuk n masukan & 25 & CPMK0607\\\garis
A2 & Tulis program yang memakai \texttt{break} untuk berhenti saat data dicari ditemukan & 25 & CPMK0607\\\garis
B1 & Tentukan keluaran perulangan bersarang yang diberikan & 20 & CPMK0608\\\garis
B2 & Turunkan rumus banyak isi per baris dari sebuah pola dan jelaskan & 20 & CPMK0608\\\garis
B3 & Jelaskan perbedaan efek \texttt{break} dan \texttt{continue} pada potongan kode & 10 & CPMK0608\\
\garisdasar
\end{tabular}}\par\vspace{6pt}
\Bagian{Penilaian}{Tugas 6: Pattern Generator}
Kumpulkan sebelum Pertemuan 7 lewat kanal pengumpulan yang diumumkan dosen.

\Sub{Bagian A: program (50 poin)}
Buat program Pattern Generator dengan ketentuan berikut.
\begin{enumerate}
\item Menu pola dengan \texttt{switch}: 1 segitiga kiri, 2 segitiga kanan, 3 piramida, 4 berlian, 5 persegi berongga, 0 keluar
\item Baca jumlah baris n
\item Tampilkan pola sesuai pilihan
\item Program berulang sampai pengguna memilih 0
\item Validasi: n harus 1 sampai 20, pilihan harus 0 sampai 5
\end{enumerate}
\begin{Contoh}[Bonus, tidak wajib]
Tambahkan pola huruf inisial nama kalian, atau pola kreasi sendiri.
\end{Contoh}
\Sub{Bagian B: penjelasan (50 poin)}
Komentar maksimal 150 kata dan tabel baris, spasi, dan isi untuk pola piramida dan berlian dengan n = 4, beserta rumus yang diturunkan dari tabel itu. Jelaskan satu kesalahan batas perulangan yang kalian temui.

\Sub{Rubrik Tugas 6}
\par\vspace{4pt}\noindent{\fontsize{8.5pt}{11.5pt}\selectfont
\begin{tabular}{@{}>{\raggedright\arraybackslash}p{0.190\textwidth}>{\raggedright\arraybackslash}p{0.070\textwidth}>{\raggedright\arraybackslash}p{0.200\textwidth}>{\raggedright\arraybackslash}p{0.180\textwidth}>{\raggedright\arraybackslash}p{0.170\textwidth}>{\raggedright\arraybackslash}p{0.170\textwidth}@{}}
\kepala{Kriteria} & \kepala{Poin} & \kepala{Sangat baik 85–100} & \kepala{Baik 70–84} & \kepala{Cukup 55–69} & \kepala{Kurang <55}\\
A1 Program berjalan & 20 & tanpa error dan peringatan & ada peringatan & perlu satu perbaikan & tidak terkompilasi\\\garis
A2 Menu dan validasi & 15 & lima pola, menu berulang, validasi & satu pola kurang & dua pola kurang & menu tidak berulang\\\garis
A3 Ketepatan bentuk & 15 & semua pola tepat untuk n = 1, 4, 7 & satu pola meleset & dua meleset & sebagian besar meleset\\\garis
B1 Tabel dan rumus & 30 & dua tabel dan rumus tepat & satu rumus kurang tepat & tabel tanpa rumus & tidak ada atau disalin\\\garis
B2 Cerita error dan pengujian & 20 & pesan, sebab, perbaikan, uji & pesan tidak dikutip & hanya perbaikan & tidak ada\\
\garisdasar
\end{tabular}}\par\vspace{6pt}
Nilai kriteria = poin × skor level. Contoh: A1 di level Baik dengan skor 80 memberi 20 × 0,80 = 16 poin.

\Bagian{Penutup}{Bacaan lanjut dan persiapan minggu depan}
\begin{enumerate}
\item Deitel, P. dan Deitel, H. C++ How to Program, edisi ke-10, Bab 5.
\item learncpp.com, Bab 8 (break, continue, nested loops).
\item Gaddis, T. Starting Out with C++, Bab 5.
\item Slide \texttt{01 Slide/P6 Perulangan Bagian 2 dan Pola.pdf} dan hands-on di \texttt{02 Hands-on/P6 - Perulangan 2 dan Pola - Hands-on/}.
\end{enumerate}
\begin{Penting}[Minggu depan]
Pertemuan 7 adalah review dan simulasi UTS untuk materi Pertemuan 1 sampai 6. Bawa catatan dan daftar hal yang masih membingungkan.
\end{Penting}
\end{document}
